%----------------------------------------------------------------------------------------
% CV content -- English and Chinese side by side. Edit this file for both CVs at once.
%
%   \EN{...}  text printed only in the English CV (cv.tex)
%   \ZH{...}  text printed only in the Chinese CV (cv_chinese.tex)
%   Text outside \EN and \ZH (dates, names, software) is printed in both.
%
% To add an item, copy one of the same kind and change it:
%   section    \begin{cvsection}{English title}{中文标题} ... \end{cvsection}
%   position   \cvposition{title}{institution}{dates}{small second line}{place}
%   school     \cvschool{school}{place}{dates}{degree lines}
%   project    \begin{rSubsubsection}{name}{dates}{}{} \item ... \end{rSubsubsection}
%   bullet     \item \EN{English text}
%                    \ZH{中文内容}
% The dropped language also drops the space or line break after it, so \EN{...} and
% \ZH{...} may sit on one line or on two. Do not put & (table columns) inside them.
% Entries of one section are separated by a blank line; list them newest first.
%
% Publications and talks are not here: add them to publications.bib / talks.bib.
%----------------------------------------------------------------------------------------

\nocite{*}

%----------------------------------------------------------------------------------------
%	SUMMARY
%----------------------------------------------------------------------------------------

\begin{cvsection}{Summary}{个人简介}
\EN{Ph.D.\ in Physics with a strong background in astroparticle physics, data analysis,
instrumentation, and collaborative team management. Aspiring to an academic career leading a
research program in astroparticle physics that bridges physics analysis methods with novel
instrumentation development for next-generation experiments.}
\ZH{美国东北大学物理学博士，在天体粒子物理、数据分析、探测器仪器研制及团队协作管理方面具有扎实的背景，参与NASA通用反粒子谱仪实验GAPS、伽马射线与反粒子探测器实验GRAMS等国际合作项目，并任GRAMS原型机项目科学。致力于在学术界建立并领导天体粒子物理研究方向，将物理分析方法与新型探测器研制相结合，服务于下一代实验。}
\end{cvsection}

%----------------------------------------------------------------------------------------
%	POSITIONS
%----------------------------------------------------------------------------------------

\begin{cvsection}{Academic Positions}{工作访学经历}
\cvposition
  {\EN{Postdoctoral Researcher}\ZH{博士后研究员}}
  {\EN{Tsung-Dao Lee Institute, Shanghai Jiao Tong University}\ZH{上海交通大学李政道研究所}}
  {\EN{Jul 2026 -- Present}\ZH{2026年7月 -- 至今}}
  {\EN{State Key Laboratory of Dark Matter Physics; Key Laboratory for Particle Astrophysics and Cosmology (MoE); \\
       Shanghai Key Laboratory for Particle Physics and Cosmology}%
   \ZH{暗物质物理全国重点实验室；粒子天体物理与宇宙学教育部重点实验室；上海市粒子物理与宇宙学重点实验室}}
  {\EN{Shanghai, China}\ZH{中国上海}}

\cvposition
  {\EN{Part-time Staff Associate}\ZH{兼职研究人员}}
  {\EN{Columbia Astrophysics Laboratory, Columbia University}\ZH{哥伦比亚大学哥伦比亚天体物理实验室}}
  {\EN{Jul 2023 -- Jun 2024}\ZH{2023年7月 -- 2024年6月}}
  {\EN{Field: Physics}\ZH{Part-time Staff Associate；专业方向：物理学}}
  {\EN{New York, NY, USA}\ZH{美国纽约}}
\end{cvsection}

%----------------------------------------------------------------------------------------
%	EDUCATION
%----------------------------------------------------------------------------------------

\begin{cvsection}{Education}{教育背景}
\cvschool
  {\EN{Northeastern University}\ZH{美国东北大学（Northeastern University）}}
  {\EN{Boston, MA, USA}\ZH{美国马萨诸塞州波士顿}}
  {\EN{2018 -- 2026}\ZH{2018 -- 2026}}
  {\EN{Master and Ph.D.\ in Physics}\ZH{物理学硕士及博士} \hfill GPA 3.5/4.0}

\cvschool
  {\EN{Sun Yat-sen University}\ZH{中山大学}}
  {\EN{Guangzhou, Guangdong, China}\ZH{中国广东广州}}
  {2014 -- 2018}
  {\EN{B.S.\ in Physics}\ZH{物理学学士} \hfill GPA 3.2/4.0}
\end{cvsection}

%----------------------------------------------------------------------------------------
%	RESEARCH EXPERIENCE
%----------------------------------------------------------------------------------------

\begin{cvsection}{Research Experience}{科研经历}

\begin{rSubsubsectionTitle}{\EN{Graduate Researcher, Northeastern University}\ZH{研究生研究员，美国东北大学}}{\EN{2018 -- 2026}\ZH{2018 -- 2026}}{}{}
\item

  \begin{rSubsubsection}{\EN{Gamma-Ray and AntiMatter Survey (GRAMS), NASA APRA Project}\ZH{伽马射线与反粒子探测器实验（GRAMS），NASA APRA项目}}{\EN{Mar 2021 -- Present}\ZH{2021年3月 -- 至今}}{}{}
  \item \textit{\EN{Leadership: Subsystem Integration Coordinator and Project Scientist, GRAMS prototype flight (pGRAMS)}\ZH{领导职务：GRAMS原型气球飞行（pGRAMS）子系统集成协调人，项目科学家}}
  \item \EN{Leading collaborative efforts in preparation for the pGRAMS launch.}
        \ZH{牵头协调pGRAMS发射的各项准备工作。}
  \item \EN{Designed and fabricated the pGRAMS time projection chamber (TPC) charge-sensitive amplifiers.}
        \ZH{设计并制作pGRAMS时间投影室（TPC）的电荷灵敏放大器。}
  \item \EN{Designed the mechanical CAD of the pGRAMS TPC detector.}
        \ZH{完成pGRAMS TPC探测器的机械CAD设计。}
  \item \EN{Fabricated the TPC and set up the chamber environment for liquid argon TPC operations.}
        \ZH{制作TPC，并搭建液氩TPC运行所需的腔体环境。}
  \item \EN{Simulated the pGRAMS TPC static electric field.}
        \ZH{模拟pGRAMS TPC的静电场分布。}
  \item \EN{Developed Geant4 simulations for GRAMS and the GRAMS antideuteron and antihelium-3 sensitivity.}
        \ZH{开发GRAMS的Geant4模拟，并完成GRAMS反氘及反氦-3灵敏度研究。}
  \end{rSubsubsection}

  \begin{rSubsubsection}{\EN{General AntiParticle Spectrometer (GAPS), NASA APRA Project}\ZH{通用反粒子谱仪实验（GAPS），NASA APRA项目}}{\EN{Mar 2021 -- Present}\ZH{2021年3月 -- 至今}}{}{}
  \item \EN{Assembled the GAPS functional prototype (GFP) at MIT Bates Lab, improvised the GAPS thermal system, and built a simulation model for the GFP.}
        \ZH{在麻省理工学院Bates实验室组装GAPS功能原型机（GFP），改进GAPS热控系统，并建立GFP模拟模型。}
  \item \EN{Integrated the GAPS payload at MIT Bates Lab, UCB Space Sciences Lab, and Columbia Nevis Labs; assembled the tracker, time-of-flight, and cooling systems.}
        \ZH{在麻省理工学院Bates实验室、加州大学伯克利分校空间科学实验室及哥伦比亚大学Nevis实验室完成GAPS载荷集成；组装径迹探测器、飞行时间系统及冷却系统。}
  \item \EN{Developed the ground cooling system that cools all GAPS detectors to $-40\,^\circ$C.}
        \ZH{研制GAPS地面冷却系统，将全部探测器冷却至$-40\,^\circ$C。}
  \item \EN{Developed software for detector performance visualization.}
        \ZH{开发探测器性能可视化软件。}
  \item \EN{Developed the GAPS atmospheric unfolding model and analysis.}
        \ZH{建立GAPS大气效应反演（unfolding）模型并完成相应分析。}
  \end{rSubsubsection}

  \begin{rSubsubsection}{\EN{Two-Photon Excitation Microscopy}\ZH{双光子激发显微成像}}{\EN{Sept 2018 -- Feb 2021}\ZH{2018年9月 -- 2021年2月}}{}{}
  \item \EN{Designed an FPGA-based data acquisition system (DAQ) and data down-sampling algorithm using StellarIP and Vivado, and calibrated 32-channel firmware for the FPGA output; used as the back end of two-photon excitation microscopy scanning for cancer.}
        \ZH{基于StellarIP与Vivado设计FPGA数据采集系统（DAQ）及数据降采样算法，并标定FPGA输出的32通道固件；该系统用作癌症双光子激发显微扫描成像的后端。}
  \end{rSubsubsection}

\end{rSubsubsectionTitle}

\begin{rSubsubsectionTitle}{\EN{Undergraduate Researcher, Sun Yat-sen University}\ZH{本科生科研，中山大学}}{2014 -- 2018}{}{}
\item

  \begin{rSubsubsection}{\EN{Heavy-Flavor Meson Decay Simulation for the PHENIX Experiment}\ZH{PHENIX实验重味介子衰变模拟}}{}{}{}
  \item \EN{Simulated B and D mesons passing through four detector planes, and used the asymmetry of muons detected on the planes to determine the initial meson ratio.}
        \ZH{模拟B介子与D介子穿过四层探测器平面的过程，利用各平面探测到的$\mu$子不对称性确定初始介子比例。}
  \end{rSubsubsection}

  \begin{rSubsubsection}{\EN{Galaxies, Cosmology, and Structure Formation Simulation}\ZH{星系、宇宙学与结构形成模拟}}{}{}{}
  \item \EN{Simulated galaxy formation using GADGET-2.}
        \ZH{使用GADGET-2模拟星系形成。}
  \end{rSubsubsection}

\end{rSubsubsectionTitle}

\end{cvsection}

%----------------------------------------------------------------------------------------
%	PUBLICATIONS  (from publications.bib, split by `keywords`)
%----------------------------------------------------------------------------------------

\begin{cvsection}{Publications}{论文发表}
\newrefcontext[labelprefix=L]
\printbibliography[keyword=lead,    heading=cvsub, title={\EN{Lead Author}\ZH{第一作者}}, resetnumbers=true]
\newrefcontext[labelprefix=C]
\printbibliography[keyword=contrib, heading=cvsub, title={\EN{Contributing Author}\ZH{合作作者}}, resetnumbers=true]
\end{cvsection}

%----------------------------------------------------------------------------------------
%	TALKS  (from talks.bib)
%----------------------------------------------------------------------------------------

\begin{cvsection}{Talks and Presentations}{学术报告}
\newrefcontext[labelprefix=T]
\printbibliography[type=misc, heading=cvnone, resetnumbers=true]
\end{cvsection}

%----------------------------------------------------------------------------------------
%	TEACHING
%----------------------------------------------------------------------------------------

\begin{cvsection}{Teaching Experience}{教学经历}

\begin{rSubsection}{\EN{Teaching Assistant Fellowship, Northeastern University}\ZH{助教奖学金（Teaching Assistant Fellowship），美国东北大学}}{2020 -- 2021}{}{}
\item \EN{Teaching Assistant in the Physics Department for two years: taught the Introductory Physics Laboratory (IPL), led the undergraduate physics Interactive Learning Sessions (ILS), and served as grader and course TA for Modern Physics.}
      \ZH{在物理系担任助教两年：讲授普通物理实验课（IPL），主持本科物理互动学习课（ILS），并担任近代物理课程助教及阅卷人。}
\end{rSubsection}

\begin{rSubsection}{\EN{Teaching Assistant, Molecular Imaging}\ZH{分子影像课程助教}}{\EN{July 2019}\ZH{2019年7月}}{}{}
\item \EN{TA for Prof.\ Craig S.~Levin (Department of Radiology, Stanford). Lectured on Radiology.}
      \ZH{担任Craig S.~Levin教授（斯坦福大学放射学系）的助教，讲授生物发光/荧光成像及双光子显微成像。}
\end{rSubsection}

\begin{rSubsection}{\EN{Teaching Assistant, Cosmology and Structure Formation}\ZH{宇宙学与结构形成课程助教}}{\EN{July -- Aug 2018}\ZH{2018年7月 -- 8月}}{}{}
\item \EN{TA for Prof.\ Mark Vogelsberger (MIT Kavli Institute for Astrophysics and Space Research). Lectured on calculus and introductory Python programming.}
      \ZH{担任Mark Vogelsberger教授（麻省理工学院Kavli天体物理与空间研究所）的助教，讲授微积分及Python编程入门。}
\end{rSubsection}

\end{cvsection}

%----------------------------------------------------------------------------------------
%	HONORS  (one line each: text \hfill {\em year} \\ -- no \\ after the last one)
%----------------------------------------------------------------------------------------

\begin{cvsection}{Honors}{荣誉奖项}
\EN{SGA Scholarship, Northeastern University}\ZH{东北大学全额奖学金} \hfill {\em 2019-2026}\\
\EN{``Best Talk,'' Introduction to Physics Research Seminar, Northeastern University}\ZH{美国东北大学物理研究入门研讨会“最佳报告奖”（Best Talk）} \hfill {\em 2022} \\
\EN{Lawrence Award for Excellence in Teaching, Northeastern University}\ZH{美国东北大学Lawrence优秀教学奖（Lawrence Award for Excellence in Teaching）} \hfill {\em 2020} \\
\EN{Excellent Student Scholarship for academic excellence, Sun Yat-sen University}\ZH{中山大学优秀学生奖学金} \hfill {\em 2014}
\end{cvsection}

%----------------------------------------------------------------------------------------
%	SKILLS
%----------------------------------------------------------------------------------------

\begin{cvsection}{Skills}{专业技能}
\begin{tabular}{@{} >{\bfseries}l @{\hspace{2ex}} p{0.8\textwidth} @{}}
\EN{Computing}\ZH{计算机}   & \cvlabel{Programming:}{编程：}Python, Geant4, ROOT, Bash, MATLAB, Unix, FPGA, Git, MySQL, \LaTeX \\
                            & \cvlabel{Visualization:}{可视化：}Matplotlib, Grafana, HepRApp, MeshLab, Femtet, VNC, Adobe Photoshop, Bambu Studio\EN{ (3D printing)}\ZH{（3D打印）} \\
                            & \cvlabel{Engineering:}{工程软件：}SolidWorks, Autodesk Fusion 360, Autodesk Eagle / EasyEDA \\
                            & \cvlabel{Machine learning:}{机器学习：}PyTorch \\[3pt]
\EN{Laboratory}\ZH{实验技能} & \cvlabel{Onsite:}{现场操作：}\EN{compressor operation and maintenance, piping, cryogenic (LAr, LN$_2$) system operation, crane and gantry operation, strapping and roping}\ZH{压缩机操作与维护、管路安装、低温（液氩、液氮）系统操作、吊车与龙门架操作、捆扎与绳索作业} \\
                            & \cvlabel{Physics lab:}{物理实验：}\EN{oscilloscope, lock-in amplifier, PCB population}\ZH{示波器、锁相放大器、PCB焊接组装} \\
\end{tabular}
\end{cvsection}
